Hem situat una partícula puntual amb una càrrega $q = 10\ \mu\mathrm{C}$ i dues partícules puntuals amb una càrrega $-q$ als vèrtexs d'un quadrat de costat $a = 1,50\ \mathrm{cm}$ tal com s'indica en la figura.

\figura[0.25]{fig1.png}

\begin{apartat}{1}
Quin és el valor de la càrrega puntual $q_0$ situada al quart vèrtex si la força elèctrica sobre la càrrega $q$ és nul·la?
\end{apartat}

\begin{apartat}{1}
Quin treball haurem de fer per a portar una càrrega puntual de $0,50\ \mu\mathrm{C}$ des d'una distància molt gran fins al centre del quadrat?
\end{apartat}

\blocetiquetat{Dada:}{$k = \dfrac{1}{4\pi\varepsilon_0} = 8,99\times 10^{9}\ \mathrm{N\,m^2\,C^{-2}}$.}
\nota{Suposeu que les velocitats inicial i final de la càrrega que portem fins al centre del quadrat són nul·les.}
